% !TeX program = lualatex
% !TeX encoding = utf8
% !TeX spellcheck = uk_UA
% !TeX root = ../../main.tex

%% --------------------------------------------------------
%% Chapter file: Instability.tex
%% --------------------------------------------------------

\graphicspath{{\currfilebase/Pictures}}

\chapter{Загальний принцип нестійкості плазми}
\label{sec:plasma-instability}\hypertarget{plasma-instability}{}

%% ========================================================
\section{Що означає нестійкість?}
%% ========================================================

До цього моменту ми розглядали рівноважні стани плазми та різні
рівняння, якими можна описувати її рух. Але саме по собі знаходження
рівноваги ще не означає, що така конфігурація може існувати тривалий час.

Уявімо, що плазма перебуває в деякому стані рівноваги. До неї неминуче
вносяться малі збурення: тепловий рух частинок, флуктуації густини,
електричного та магнітного полів, похибки зовнішнього поля. Питання
стійкості полягає в тому, що відбувається з таким малим відхиленням.

Якщо після відхилення система повертається до початкового стану,
рівновага стійка. Якщо збурення залишається приблизно сталим, маємо
нейтральну стійкість. Якщо ж навіть дуже мале збурення з часом
зростає, рівновага нестійка.

\begin{figure}[ht]
\centering
\fbox{%
  \parbox[c][55mm][c]{0.82\textwidth}{%
    \centering
    \textbf{ЗАГЛУШКА ДЛЯ РИСУНКА}\\[2mm]
    Три якісні графіки $|\delta X(t)|$: \\[1mm]
    \textit{стійкість} --- збурення затухає; \qquad
    \textit{нейтральність} --- збурення стало; \qquad
    \textit{нестійкість} --- збурення зростає.\\[3mm]
    По горизонталі --- час $t$, по вертикалі --- амплітуда збурення.
  }%
}
\caption{Три можливі типи еволюції малого збурення рівноваги.}
\label{fig:instability-three-regimes}
\end{figure}

Це визначення важливе тим, що воно не прив'язане до конкретної моделі.
Одна й та сама ідея працює для механічної системи, рідини, плазми,
електромагнітного поля або кінетичного розподілу частинок.

Розглянемо рівноважний стан $X_0$ і мале збурення $\delta X$. Тоді
записуємо

\begin{equation}
X=X_0+\delta X,
\label{eq:instability-perturbation}
\end{equation}

де припускаємо, що $|\delta X|\ll |X_0|$. Саме ця нерівність дозволяє
спочатку розглядати збурення незалежно від його нелінійної взаємодії
із самою рівновагою.

%% ========================================================
\section{Лінеаризація: як знайти долю малого збурення}
%% ========================================================

Нехай рівняння плазми можна символічно записати як

\begin{equation}
\mathcal{F}(X)=0.
\label{eq:instability-equilibrium-equation}
\end{equation}

Підставимо в нього представлення \eqref{eq:instability-perturbation}.
Розкладаючи рівняння за степенями малого збурення, одержуємо

\begin{equation}
\mathcal{F}(X_0)
+
\left.
\frac{\partial\mathcal{F}}{\partial X}
\right|_{X_0}\delta X
+\mathcal{O}(\delta X^2)=0.
\label{eq:instability-linearization}
\end{equation}

Перший доданок у \eqref{eq:instability-linearization} зникає, оскільки
$X_0$ є рівноважним станом. Залишається рівняння першого порядку за
збуренням:

\begin{equation}
\mathcal{L}\,\delta X=0,
\label{eq:instability-linear-equation}
\end{equation}

де $\mathcal{L}$ --- лінійний оператор, визначений властивостями
рівноваги.

Це і є головний технічний крок теорії стійкості. Замість повного
нелінійного рівняння ми спочатку питаємо значно простіше:

\begin{quote}
\textit{Що зробить система з нескінченно малим відхиленням від рівноваги?}
\end{quote}

Якщо воно зростає, рівновага нестійка вже на лінійному рівні.

%% ========================================================
\section{Хвиля, що може рости}
%% ========================================================

Для однорідної рівноваги зручно розглядати окрему просторово-часову
гармоніку:

\begin{equation}
\delta X
\propto
\exp\!\left[i(\vect k\cdot\vect r-\omega t)\right].
\label{eq:instability-wave}
\end{equation}

Тут $\vect k$ визначає просторовий масштаб збурення, а $\omega$ ---
його частоту. Важливо, що в загальному випадку $\omega$ може бути
комплексною:

\begin{equation}
\omega=\omega_r+i\gamma.
\label{eq:instability-complex-frequency}
\end{equation}

Підставивши \eqref{eq:instability-complex-frequency} у
\eqref{eq:instability-wave}, одержуємо

\begin{equation}
\delta X
\propto
\exp(i\vect k\cdot\vect r-i\omega_r t)\,
\exp(\gamma t).
\label{eq:instability-amplitude}
\end{equation}

Тому знак $\gamma$ безпосередньо визначає долю збурення:

\begin{equation}
\gamma>0 \quad\Rightarrow\quad \text{зростання},
\qquad
\gamma<0 \quad\Rightarrow\quad \text{затухання},
\qquad
\gamma=0 \quad\Rightarrow\quad \text{нейтральність}.
\label{eq:instability-criterion}
\end{equation}

Отже, критерій лінійної нестійкості можна сформулювати дуже коротко:
потрібна наявність розв'язку з

\begin{equation}
\operatorname{Im}\omega>0.
\label{eq:instability-growth-condition}
\end{equation}

Часова шкала розвитку нестійкості визначається величиною $\gamma$:

\begin{equation}
\tau_{\rm inst}\sim\frac{1}{\gamma}.
\label{eq:instability-time}
\end{equation}

Таким чином, нестійкість --- це не просто твердження, що
``плазма поводиться хаотично''. На початковій стадії це цілком
визначений експоненційний ріст певної моди збурення.

\begin{figure}[ht]
\centering
\fbox{%
  \parbox[c][58mm][c]{0.82\textwidth}{%
    \centering
    \textbf{ЗАГЛУШКА ДЛЯ РИСУНКА}\\[2mm]
    Просторово-часова картина нестійкої моди.\\
    Показати хвилю $\delta X$ вздовж координати $x$ у кілька
    послідовних моментів часу.\\[2mm]
    Амплітуда хвилі з часом зростає, а її просторовий масштаб
    залишається приблизно сталим.
  }%
}
\caption{Нестійка хвиля: просторова структура моди та її зростання
з часом.}
\label{fig:instability-growing-mode}
\end{figure}

%% ========================================================
\section{Дисперсійне рівняння}
%% ========================================================

Після підстановки \eqref{eq:instability-wave} лінеаризовані рівняння
перетворюються на систему алгебраїчних рівнянь для амплітуд різних
величин. Її можна символічно записати у вигляді

\begin{equation}
\mathsf{D}(\omega,\vect k)\,\vect A=0,
\label{eq:instability-matrix}
\end{equation}

де $\vect A$ містить амплітуди збурених величин, а
$\mathsf{D}$ --- матриця, яка визначається моделлю плазми та вибраною
рівновагою.

Нетривіальний розв'язок \eqref{eq:instability-matrix} існує лише тоді,
коли

\begin{equation}
\det\mathsf{D}(\omega,\vect k)=0.
\label{eq:instability-dispersion}
\end{equation}

Це дисперсійне рівняння. Воно є одним із центральних інструментів
теорії плазмових нестійкостей.

Для кожного $\vect k$ воно визначає можливі значення $\omega$.
Якщо серед них є корінь, що задовольняє умову
\eqref{eq:instability-growth-condition}, відповідна мода нестійка.

Тут виникає важливий зв'язок із попередніми розділами. У теорії хвиль
ми вже шукали дисперсійні співвідношення. Але тепер нас цікавить не
лише залежність дійсної частоти від хвильового числа. Насамперед треба
з'ясувати, чи має частота уявну частину і що означає її знак.

\begin{figure}[ht]
\centering
\fbox{%
  \parbox[c][58mm][c]{0.82\textwidth}{%
    \centering
    \textbf{ЗАГЛУШКА ДЛЯ РИСУНКА}\\[2mm]
    Діаграма $\omega(k)$ для кількох гілок.\\
    Окремо показати область, де $\omega$ дійсна, та область
    нестійкості, де $\omega=\omega_r+i\gamma$ і $\gamma>0$.\\[2mm]
    Бажано підписати: ``хвиля'', ``затухаюча мода'',
    ``нестійка мода''.
  }%
}
\caption{Дисперсійне співвідношення як спосіб визначення спектра
плазмових мод та їхньої стійкості.}
\label{fig:instability-dispersion}
\end{figure}

%% ========================================================
\section{Звідки береться енергія нестійкості?}
%% ========================================================

Формальна умова \eqref{eq:instability-growth-condition} говорить,
\emph{що} відбувається, але не пояснює, \emph{чому} збурення росте.
Фізично нестійкість означає, що система має джерело вільної енергії,
яку збурення може перетворювати на енергію руху, електричних або
магнітних полів.

У плазмі такою вільною енергією можуть бути:

\begin{itemize}
\item градієнти густини та температури;
\item запас енергії магнітного поля;
\item енергія електричного струму;
\item анізотропія розподілу швидкостей;
\item енергія пучка заряджених частинок;
\item неоднорідність магнітної конфігурації.
\end{itemize}

Тому нестійкість не ``виникає з нічого''. Вона є механізмом
перетворення енергії нерівноважного стану на зростаюче збурення.

\begin{figure}[ht]
\centering
\fbox{%
  \parbox[c][62mm][c]{0.82\textwidth}{%
    \centering
    \textbf{ЗАГЛУШКА ДЛЯ РИСУНКА}\\[2mm]
    Схема енергетичного механізму нестійкості:\\[1mm]
    \textit{вільна енергія рівноваги}
    $\longrightarrow$
    \textit{збурення}
    $\longrightarrow$
    \textit{енергія руху / поля}.\\[3mm]
    Навести 3--4 джерела вільної енергії:
    градієнт тиску, струм, анізотропія, пучок частинок.
  }%
}
\caption{Фізичне походження плазмової нестійкості: збурення отримує
енергію з нерівноважного стану плазми.}
\label{fig:instability-free-energy}
\end{figure}

Саме тому стабільність не можна визначити лише за величиною
зовнішнього магнітного поля. Сильне поле може змінювати траєкторії
частинок, але якщо в системі залишається доступне джерело вільної
енергії, можуть існувати нестійкі моди.

%% ========================================================
\section{Рівновага та стійкість у магнітному полі}
%% ========================================================

Для магнітного утримання особливо важливо розрізняти два питання:

\begin{enumerate}
\item чи може існувати рівновага сил;
\item чи є ця рівновага стійкою до малих збурень.
\end{enumerate}

У простому однорідному МГД-описі рівновага плазми задовольняє
умову

\begin{equation}
-\nabla p+\frac{1}{c}[\vect j\times\Bfield]=0.
\label{eq:instability-mhd-equilibrium}
\end{equation}

Рівняння \eqref{eq:instability-mhd-equilibrium} лише говорить, що
сила градієнта тиску компенсується магнітною силою. Воно ще нічого
не говорить про те, що стан буде зберігатися після малого зміщення.

Щоб зрозуміти це фізично, розглянемо магнітну силу. Її зручно
розкласти на дві частини:

\begin{equation}
\frac{1}{c}[\vect j\times\Bfield]
=
-\nabla\frac{B^2}{8\pi}
+
\frac{1}{4\pi}(\Bfield\cdot\nabla)\Bfield.
\label{eq:instability-magnetic-force}
\end{equation}

Перша частина поводиться як градієнт магнітного тиску

\begin{equation}
p_B=\frac{B^2}{8\pi},
\label{eq:instability-magnetic-pressure}
\end{equation}

а друга пов'язана з натягом викривлених силових ліній.

Отже, магнітне поле не просто ``тисне'' на плазму. Його силові
лінії можуть діяти подібно до натягнутих еластичних ниток, намагаючись
зменшити свою кривину.

\begin{figure}[ht]
\centering
\fbox{%
  \parbox[c][65mm][c]{0.82\textwidth}{%
    \centering
    \textbf{ЗАГЛУШКА ДЛЯ РИСУНКА}\\[2mm]
    Магнітна силова лінія до та після малого викривлення.\\[2mm]
    Показати дві дії магнітного поля:
    \begin{itemize}
      \item градієнт $B^2/8\pi$ як ``магнітний тиск'';
      \item натяг силової лінії як відновлювальну силу.
    \end{itemize}
    На рисунку бажано буквально позначити
    ``тиск'' та ``натяг''.
  }%
}
\caption{Дві складові магнітної сили: магнітний тиск і натяг силових
ліній.}
\label{fig:magnetic-pressure-tension}
\end{figure}

Ця картина стане особливо важливою у наступному розділі, де
розглядатимуться конкретні конфігурації магнітного утримання.

%% ========================================================
\section{Перетяжка як найпростіший приклад нестійкості}
%% ========================================================

Уявімо плазмовий струм, який створює власне магнітне поле. Якщо
плазмовий канал має невелику локальну перетяжку, магнітна сила може
зробити цю перетяжку ще сильнішою. Вузька ділянка стискається,
поле там посилюється, а це змінює розподіл сил так, що початкове
збурення не обов'язково буде компенсоване.

Це принципово відрізняється від звичайної стійкої системи. Для
стійкої рівноваги початкове зміщення породжує силу, яка повертає
систему назад. Для нестійкої --- виникає сила в напрямку самого
зміщення.

\begin{figure}[ht]
\centering
\fbox{%
  \parbox[c][68mm][c]{0.82\textwidth}{%
    \centering
    \textbf{ЗАГЛУШКА ДЛЯ РИСУНКА}\\[2mm]
    Схема перетяжкової нестійкості плазмового каналу.\\[2mm]
    Ліворуч: початкова майже циліндрична конфігурація.\\
    По центру: мала перетяжка як збурення.\\
    Праворуч: перетяжка посилюється.\\[2mm]
    Показати напрямки струму, магнітного поля та сили,
    яка підсилює збурення.
  }%
}
\caption{Фізична картина перетяжкової нестійкості: мале збурення
може породжувати силу, яка його підсилює.}
\label{fig:pinch-instability}
\end{figure}

Саме такі міркування лежать в основі аналізу пінч-конфігурацій,
який буде проведено в \S\ref{sec:magnetic-confinement}.

%% ========================================================
\section{Лінійна та нелінійна стадії}
%% ========================================================

Поки виконується умова малості збурення в \eqref{eq:instability-perturbation},
можна користуватися лінеаризованим рівнянням \eqref{eq:instability-linear-equation}.
Тоді кожна мода розвивається незалежно від інших, а її амплітуда
для нестійкого випадку зростає як

\begin{equation}
|\delta X(t)|
\sim
|\delta X(0)|e^{\gamma t}.
\label{eq:instability-linear-growth}
\end{equation}

Але експоненційне зростання не може тривати нескінченно. Рано чи
пізно збурення стає порівнянним із величиною рівноважних величин,
і члени вищих порядків, якими ми нехтували при лінеаризації,
стають важливими.

Починається нелінійна стадія. На ній можливі насичення нестійкості,
перерозподіл енергії між модами, утворення структур, турбулентність
або руйнування початкової конфігурації.

\begin{figure}[ht]
\centering
\fbox{%
  \parbox[c][58mm][c]{0.82\textwidth}{%
    \centering
    \textbf{ЗАГЛУШКА ДЛЯ РИСУНКА}\\[2mm]
    Часова еволюція нестійкості: \\[1mm]
    \textit{малий шум}
    $\rightarrow$
    \textit{лінійний експоненційний ріст}
    $\rightarrow$
    \textit{нелінійна стадія}
    $\rightarrow$
    \textit{насичення / перебудова}.\\[2mm]
    Виділити вертикальною або пунктирною лінією межу
    застосовності лінеаризації.
  }%
}
\caption{Типова еволюція нестійкості від малого збурення до
нелінійної стадії.}
\label{fig:instability-linear-nonlinear}
\end{figure}

Це розрізнення буде важливим далі: лінійна теорія визначає, \emph{чи}
є рівновага нестійкою і які моди ростуть найшвидше, але не завжди
визначає кінцевий стан системи.

%% ========================================================
\section{Від загального принципу до конкретних нестійкостей}
%% ========================================================

Отже, аналіз плазмової нестійкості має загальну послідовність:

\begin{enumerate}
\item задається рівноважний стан;
\item до нього додається мале збурення;
\item рівняння лінеаризуються;
\item шукаються моди у вигляді \eqref{eq:instability-wave};
\item з отриманого дисперсійного рівняння \eqref{eq:instability-dispersion}
      визначається $\omega(\vect k)$;
\item знак $\operatorname{Im}\omega$ визначає стійкість моди;
\item для нестійких мод визначається швидкість їхнього росту
      за \eqref{eq:instability-time}.
\end{enumerate}

Фізично ж треба поставити ще одне питання: \emph{звідки ця мода
бере енергію?} Саме відповідь на нього відрізняє різні класи
нестійкостей.

У магнітному утриманні такими джерелами можуть бути струми,
градієнти тиску, кривина та неоднорідність магнітного поля. У
кінетичній плазмі додатково виникають анізотропія розподілу
швидкостей, пучки частинок і резонансна взаємодія хвиль із частинками.

Таким чином, наступний крок --- не вводити ще одну абстрактну
формулу, а розглянути конкретні конфігурації та з'ясувати, які
збурення вони допускають і які фізичні механізми їх підсилюють.

Саме з цього починається аналіз магнітного утримання плазми.

%% --------------------------------------------------------
\section*{Задачі}
%% --------------------------------------------------------

\begin{problem}\label{prb:PlasmaInstabilityGrowth}
Для малого збурення, що має вигляд
$\delta X\propto\exp[i(\vect k\cdot\vect r-\omega t)]$,
знайдіть час зростання збурення, якщо
$\omega=\omega_r+i\gamma$ і
$\gamma=2\times10^5~\mathrm{s}^{-1}$.
Що зміниться, якщо $\gamma<0$?
\end{problem}

%% --------------------------------------------------------
\section*{Контрольні питання}
%% --------------------------------------------------------

\begin{Questions}
\begin{quizlist}

\item Що фізично означає стійкість, нейтральність і нестійкість
рівноваги?

\item Чому при дослідженні нестійкості спочатку можна розглядати
лише лінійні за збуренням члени?

\item Який фізичний зміст дійсної та уявної частин частоти
$\omega$?

\item Чому умова $\operatorname{Im}\omega>0$ означає зростання
збурення?

\item Чим відрізняються лінійна та нелінійна стадії розвитку
нестійкості?

\item Чому виконання рівняння рівноваги ще не гарантує стійкості
плазмової конфігурації?

\item Які дві фізичні ролі магнітного поля видно з
\eqref{eq:instability-magnetic-force}?

\item Звідки плазмова нестійкість може отримувати вільну енергію?

\end{quizlist}
\end{Questions}

